\clearpage\phantomsection\label{couverture-programme}\pdfbookmark[0]{Repérage du programme}{bm-couverture-programme}\gdef\otChapNum{0}\pagestyle{otFront}
\otPageTitre{Repérage du programme}{CM2 · Programme 2025 du cycle 3}
Cette matrice relie les objectifs d’apprentissage du programme aux
exercices du livre. Les cours complètent ces repères, notamment pour
le vocabulaire, les représentations et les méthodes de construction.\par\bigskip
Les six domaines du programme sont couverts : nombres, calcul et
résolution de problèmes ; grandeurs et mesures ; espace et géométrie ;
organisation et gestion de données et probabilités ; proportionnalité ;
initiation à la pensée informatique, travaillée dans les chapitres 5,
15 et 18 ainsi que dans plusieurs défis.\par\bigskip
\begin{attention}Les niveaux de difficulté et les défis sont des choix
d’entraînement. Ils ne constituent pas une liste d’obligations
supplémentaires au programme officiel.\end{attention}\clearpage
\otPageTitre{Repères et exercices}{Chapitres 1 à 3}\fontsize{10}{13}\selectfont
\otIntertitre{1. Les nombres jusqu’au milliard}\begin{tabularx}{\linewidth}{@{}Yl@{}}\toprule
Classes des millions et lectures & \hyperref[exo-1-1]{1.1}, \hyperref[exo-1-2]{1.2}, \hyperref[exo-1-13]{1.13}\\[2mm]
Décompositions et valeurs & \hyperref[exo-1-3]{1.3}, \hyperref[exo-1-4]{1.4}, \hyperref[exo-1-6]{1.6}, \hyperref[exo-1-12]{1.12}\\[2mm]
Comparer, encadrer, intercaler & \hyperref[exo-1-5]{1.5}, \hyperref[exo-1-8]{1.8}, \hyperref[exo-1-9]{1.9}, \hyperref[exo-1-10]{1.10}, \hyperref[exo-1-14]{1.14}, \hyperref[exo-1-15]{1.15}\\[2mm]
Demi-droite et problèmes & \hyperref[exo-1-7]{1.7}, \hyperref[exo-1-11]{1.11}\\[2mm]
\bottomrule\end{tabularx}\par\vspace{4mm}
\otIntertitre{2. Additionner, soustraire, multiplier}\begin{tabularx}{\linewidth}{@{}Yl@{}}\toprule
Techniques posées et contrôles & \hyperref[exo-2-1]{2.1}, \hyperref[exo-2-2]{2.2}, \hyperref[exo-2-5]{2.5}, \hyperref[exo-2-10]{2.10}\\[2mm]
Calcul réfléchi & \hyperref[exo-2-3]{2.3}, \hyperref[exo-2-4]{2.4}, \hyperref[exo-2-15]{2.15}\\[2mm]
Parenthèses et expressions & \hyperref[exo-2-6]{2.6}, \hyperref[exo-2-7]{2.7}, \hyperref[exo-2-8]{2.8}, \hyperref[exo-2-9]{2.9}, \hyperref[exo-2-14]{2.14}\\[2mm]
Problèmes et factures & \hyperref[exo-2-11]{2.11}, \hyperref[exo-2-12]{2.12}, \hyperref[exo-2-13]{2.13}\\[2mm]
\bottomrule\end{tabularx}\par\vspace{4mm}
\otIntertitre{3. Multiples, diviseurs et partages}\begin{tabularx}{\linewidth}{@{}Yl@{}}\toprule
Multiples et critères & \hyperref[exo-3-1]{3.1}, \hyperref[exo-3-2]{3.2}, \hyperref[exo-3-5]{3.5}, \hyperref[exo-3-8]{3.8}, \hyperref[exo-3-10]{3.10}, \hyperref[exo-3-15]{3.15}\\[2mm]
Listes complètes de diviseurs & \hyperref[exo-3-3]{3.3}, \hyperref[exo-3-4]{3.4}, \hyperref[exo-3-9]{3.9}, \hyperref[exo-3-14]{3.14}\\[2mm]
Diviseurs et multiples communs & \hyperref[exo-3-6]{3.6}, \hyperref[exo-3-7]{3.7}\\[2mm]
Partages et rendez-vous & \hyperref[exo-3-11]{3.11}, \hyperref[exo-3-12]{3.12}, \hyperref[exo-3-13]{3.13}\\[2mm]
\bottomrule\end{tabularx}\par\vspace{4mm}
\clearpage
\otPageTitre{Repères et exercices}{Chapitres 4 à 6}\fontsize{10}{13}\selectfont
\otIntertitre{4. La division : entière et décimale}\begin{tabularx}{\linewidth}{@{}Yl@{}}\toprule
Division euclidienne contrôlée & \hyperref[exo-4-1]{4.1}, \hyperref[exo-4-2]{4.2}, \hyperref[exo-4-7]{4.7}, \hyperref[exo-4-9]{4.9}\\[2mm]
Quotient décimal exact & \hyperref[exo-4-3]{4.3}, \hyperref[exo-4-4]{4.4}, \hyperref[exo-4-5]{4.5}, \hyperref[exo-4-8]{4.8}, \hyperref[exo-4-10]{4.10}\\[2mm]
Choisir le quotient utile & \hyperref[exo-4-6]{4.6}, \hyperref[exo-4-12]{4.12}, \hyperref[exo-4-15]{4.15}\\[2mm]
Partages et problèmes & \hyperref[exo-4-11]{4.11}, \hyperref[exo-4-13]{4.13}, \hyperref[exo-4-14]{4.14}\\[2mm]
\bottomrule\end{tabularx}\par\vspace{4mm}
\otIntertitre{5. Égalités, programmes et suites}\begin{tabularx}{\linewidth}{@{}Yl@{}}\toprule
Égalités et nombres cachés & \hyperref[exo-5-1]{5.1}, \hyperref[exo-5-7]{5.7}, \hyperref[exo-5-10]{5.10}, \hyperref[exo-5-15]{5.15}\\[2mm]
Programmes : exécuter et produire & \hyperref[exo-5-2]{5.2}, \hyperref[exo-5-3]{5.3}, \hyperref[exo-5-4]{5.4}, \hyperref[exo-5-8]{5.8}, \hyperref[exo-5-13]{5.13}\\[2mm]
Suites de nombres & \hyperref[exo-5-5]{5.5}, \hyperref[exo-5-6]{5.6}, \hyperref[exo-5-9]{5.9}, \hyperref[exo-5-14]{5.14}\\[2mm]
Motifs et étapes lointaines & \hyperref[exo-5-11]{5.11}, \hyperref[exo-5-12]{5.12}\\[2mm]
\bottomrule\end{tabularx}\par\vspace{4mm}
\otIntertitre{6. Les fractions, du tout au nombre}\begin{tabularx}{\linewidth}{@{}Yl@{}}\toprule
Parts et dénominateurs jusqu’à 60 & \hyperref[exo-6-1]{6.1}, \hyperref[exo-6-2]{6.2}, \hyperref[exo-6-15]{6.15}\\[2mm]
Entier plus fraction & \hyperref[exo-6-3]{6.3}, \hyperref[exo-6-4]{6.4}, \hyperref[exo-6-9]{6.9}, \hyperref[exo-6-13]{6.13}\\[2mm]
Demi-droite et repérages & \hyperref[exo-6-6]{6.6}, \hyperref[exo-6-7]{6.7}, \hyperref[exo-6-8]{6.8}\\[2mm]
Encadrements et problèmes & \hyperref[exo-6-5]{6.5}, \hyperref[exo-6-10]{6.10}, \hyperref[exo-6-11]{6.11}, \hyperref[exo-6-12]{6.12}, \hyperref[exo-6-14]{6.14}\\[2mm]
\bottomrule\end{tabularx}\par\vspace{4mm}
\clearpage
\otPageTitre{Repères et exercices}{Chapitres 7 à 9}\fontsize{10}{13}\selectfont
\otIntertitre{7. Comparer et additionner des fractions}\begin{tabularx}{\linewidth}{@{}Yl@{}}\toprule
Comparaisons directes & \hyperref[exo-7-1]{7.1}, \hyperref[exo-7-2]{7.2}, \hyperref[exo-7-5]{7.5}\\[2mm]
Écritures équivalentes & \hyperref[exo-7-3]{7.3}, \hyperref[exo-7-4]{7.4}, \hyperref[exo-7-14]{7.14}\\[2mm]
Sommes et différences compatibles & \hyperref[exo-7-6]{7.6}, \hyperref[exo-7-7]{7.7}, \hyperref[exo-7-8]{7.8}, \hyperref[exo-7-9]{7.9}, \hyperref[exo-7-10]{7.10}, \hyperref[exo-7-15]{7.15}\\[2mm]
Problèmes de parts & \hyperref[exo-7-11]{7.11}, \hyperref[exo-7-12]{7.12}, \hyperref[exo-7-13]{7.13}\\[2mm]
\bottomrule\end{tabularx}\par\vspace{4mm}
\otIntertitre{8. Prendre une fraction d’une quantité}\begin{tabularx}{\linewidth}{@{}Yl@{}}\toprule
Fraction d’une quantité & \hyperref[exo-8-1]{8.1}, \hyperref[exo-8-2]{8.2}, \hyperref[exo-8-5]{8.5}, \hyperref[exo-8-8]{8.8}\\[2mm]
Produit entier par fraction & \hyperref[exo-8-3]{8.3}, \hyperref[exo-8-4]{8.4}, \hyperref[exo-8-9]{8.9}, \hyperref[exo-8-14]{8.14}\\[2mm]
Retrouver le tout & \hyperref[exo-8-7]{8.7}, \hyperref[exo-8-15]{8.15}\\[2mm]
Problèmes à étapes & \hyperref[exo-8-6]{8.6}, \hyperref[exo-8-10]{8.10}, \hyperref[exo-8-11]{8.11}, \hyperref[exo-8-12]{8.12}, \hyperref[exo-8-13]{8.13}\\[2mm]
\bottomrule\end{tabularx}\par\vspace{4mm}
\otIntertitre{9. Les décimaux jusqu’aux millièmes}\begin{tabularx}{\linewidth}{@{}Yl@{}}\toprule
Millièmes et décompositions & \hyperref[exo-9-1]{9.1}, \hyperref[exo-9-2]{9.2}, \hyperref[exo-9-3]{9.3}, \hyperref[exo-9-4]{9.4}, \hyperref[exo-9-13]{9.13}\\[2mm]
Comparer et ranger & \hyperref[exo-9-5]{9.5}, \hyperref[exo-9-6]{9.6}, \hyperref[exo-9-11]{9.11}, \hyperref[exo-9-12]{9.12}, \hyperref[exo-9-15]{9.15}\\[2mm]
Encadrer, intercaler & \hyperref[exo-9-7]{9.7}, \hyperref[exo-9-8]{9.8}, \hyperref[exo-9-10]{9.10}\\[2mm]
Arrondis et problèmes & \hyperref[exo-9-9]{9.9}, \hyperref[exo-9-14]{9.14}\\[2mm]
\bottomrule\end{tabularx}\par\vspace{4mm}
\clearpage
\otPageTitre{Repères et exercices}{Chapitres 10 à 12}\fontsize{10}{13}\selectfont
\otIntertitre{10. Calculer avec les décimaux}\begin{tabularx}{\linewidth}{@{}Yl@{}}\toprule
Poser avec les décimaux & \hyperref[exo-10-1]{10.1}, \hyperref[exo-10-2]{10.2}, \hyperref[exo-10-9]{10.9}, \hyperref[exo-10-10]{10.10}\\[2mm]
Rangs : par 10, 100, 1 000 & \hyperref[exo-10-3]{10.3}, \hyperref[exo-10-15]{10.15}\\[2mm]
Doubles, moitiés, divisions & \hyperref[exo-10-4]{10.4}, \hyperref[exo-10-5]{10.5}, \hyperref[exo-10-6]{10.6}, \hyperref[exo-10-7]{10.7}, \hyperref[exo-10-8]{10.8}\\[2mm]
Problèmes de la vie courante & \hyperref[exo-10-11]{10.11}, \hyperref[exo-10-12]{10.12}, \hyperref[exo-10-13]{10.13}, \hyperref[exo-10-14]{10.14}\\[2mm]
\bottomrule\end{tabularx}\par\vspace{4mm}
\otIntertitre{11. Longueurs, masses, contenances}\begin{tabularx}{\linewidth}{@{}Yl@{}}\toprule
Relations et conversions & \hyperref[exo-11-1]{11.1}, \hyperref[exo-11-2]{11.2}, \hyperref[exo-11-3]{11.3}, \hyperref[exo-11-15]{11.15}\\[2mm]
Estimations et vraisemblance & \hyperref[exo-11-4]{11.4}, \hyperref[exo-11-8]{11.8}\\[2mm]
Mesures combinées & \hyperref[exo-11-5]{11.5}, \hyperref[exo-11-6]{11.6}, \hyperref[exo-11-7]{11.7}, \hyperref[exo-11-13]{11.13}\\[2mm]
Problèmes à étapes & \hyperref[exo-11-9]{11.9}, \hyperref[exo-11-10]{11.10}, \hyperref[exo-11-11]{11.11}, \hyperref[exo-11-12]{11.12}, \hyperref[exo-11-14]{11.14}\\[2mm]
\bottomrule\end{tabularx}\par\vspace{4mm}
\otIntertitre{12. Les aires : unités et formules}\begin{tabularx}{\linewidth}{@{}Yl@{}}\toprule
Unités d’aire et conversions & \hyperref[exo-12-4]{12.4}, \hyperref[exo-12-5]{12.5}, \hyperref[exo-12-8]{12.8}\\[2mm]
Formules du carré et du rectangle & \hyperref[exo-12-1]{12.1}, \hyperref[exo-12-2]{12.2}, \hyperref[exo-12-6]{12.6}, \hyperref[exo-12-10]{12.10}\\[2mm]
Figures composées & \hyperref[exo-12-7]{12.7}, \hyperref[exo-12-12]{12.12}\\[2mm]
Aires, périmètres et prix & \hyperref[exo-12-3]{12.3}, \hyperref[exo-12-9]{12.9}, \hyperref[exo-12-11]{12.11}, \hyperref[exo-12-13]{12.13}, \hyperref[exo-12-14]{12.14}, \hyperref[exo-12-15]{12.15}\\[2mm]
\bottomrule\end{tabularx}\par\vspace{4mm}
\clearpage
\otPageTitre{Repères et exercices}{Chapitres 13 à 15}\fontsize{10}{13}\selectfont
\otIntertitre{13. Les angles}\begin{tabularx}{\linewidth}{@{}Yl@{}}\toprule
Vocabulaire et comparaisons & \hyperref[exo-13-1]{13.1}, \hyperref[exo-13-2]{13.2}, \hyperref[exo-13-10]{13.10}, \hyperref[exo-13-14]{13.14}\\[2mm]
Reports et reproductions & \hyperref[exo-13-3]{13.3}, \hyperref[exo-13-15]{13.15}\\[2mm]
Sommes, doubles et moitiés & \hyperref[exo-13-6]{13.6}, \hyperref[exo-13-7]{13.7}, \hyperref[exo-13-8]{13.8}, \hyperref[exo-13-9]{13.9}\\[2mm]
L’angle droit et ses 90 degrés & \hyperref[exo-13-4]{13.4}, \hyperref[exo-13-5]{13.5}, \hyperref[exo-13-11]{13.11}, \hyperref[exo-13-12]{13.12}, \hyperref[exo-13-13]{13.13}\\[2mm]
\bottomrule\end{tabularx}\par\vspace{4mm}
\otIntertitre{14. Durées : jusqu’à la seconde}\begin{tabularx}{\linewidth}{@{}Yl@{}}\toprule
Secondes et conversions & \hyperref[exo-14-1]{14.1}, \hyperref[exo-14-2]{14.2}, \hyperref[exo-14-3]{14.3}, \hyperref[exo-14-9]{14.9}\\[2mm]
Sommes et retenues en base 60 & \hyperref[exo-14-4]{14.4}, \hyperref[exo-14-11]{14.11}, \hyperref[exo-14-15]{14.15}\\[2mm]
Instants, débuts et fins & \hyperref[exo-14-6]{14.6}, \hyperref[exo-14-7]{14.7}, \hyperref[exo-14-8]{14.8}, \hyperref[exo-14-14]{14.14}\\[2mm]
Chronos et problèmes & \hyperref[exo-14-5]{14.5}, \hyperref[exo-14-10]{14.10}, \hyperref[exo-14-12]{14.12}, \hyperref[exo-14-13]{14.13}\\[2mm]
\bottomrule\end{tabularx}\par\vspace{4mm}
\otIntertitre{15. Programmes de construction}\begin{tabularx}{\linewidth}{@{}Yl@{}}\toprule
Tracés aux instruments & \hyperref[exo-15-1]{15.1}, \hyperref[exo-15-2]{15.2}, \hyperref[exo-15-3]{15.3}\\[2mm]
Cercle et distances & \hyperref[exo-15-4]{15.4}, \hyperref[exo-15-8]{15.8}, \hyperref[exo-15-12]{15.12}\\[2mm]
Exécuter un programme & \hyperref[exo-15-5]{15.5}, \hyperref[exo-15-6]{15.6}, \hyperref[exo-15-14]{15.14}\\[2mm]
Élaborer et corriger un programme & \hyperref[exo-15-7]{15.7}, \hyperref[exo-15-9]{15.9}, \hyperref[exo-15-10]{15.10}, \hyperref[exo-15-11]{15.11}, \hyperref[exo-15-13]{15.13}, \hyperref[exo-15-15]{15.15}\\[2mm]
\bottomrule\end{tabularx}\par\vspace{4mm}
\clearpage
\otPageTitre{Repères et exercices}{Chapitres 16 à 18}\fontsize{10}{13}\selectfont
\otIntertitre{16. Polygones, du triangle à l’hexagone}\begin{tabularx}{\linewidth}{@{}Yl@{}}\toprule
Triangles et constructions & \hyperref[exo-16-2]{16.2}, \hyperref[exo-16-3]{16.3}\\[2mm]
Rectangles, losanges, trapèzes & \hyperref[exo-16-4]{16.4}, \hyperref[exo-16-5]{16.5}, \hyperref[exo-16-7]{16.7}, \hyperref[exo-16-8]{16.8}, \hyperref[exo-16-11]{16.11}, \hyperref[exo-16-13]{16.13}\\[2mm]
Pentagones et hexagones & \hyperref[exo-16-1]{16.1}, \hyperref[exo-16-9]{16.9}, \hyperref[exo-16-10]{16.10}, \hyperref[exo-16-12]{16.12}, \hyperref[exo-16-14]{16.14}\\[2mm]
Justifier aux instruments & \hyperref[exo-16-6]{16.6}, \hyperref[exo-16-15]{16.15}\\[2mm]
\bottomrule\end{tabularx}\par\vspace{4mm}
\otIntertitre{17. La symétrie, axes droits et obliques}\begin{tabularx}{\linewidth}{@{}Yl@{}}\toprule
Axes et figures usuelles & \hyperref[exo-17-1]{17.1}, \hyperref[exo-17-8]{17.8}, \hyperref[exo-17-13]{17.13}\\[2mm]
Axes verticaux et horizontaux & \hyperref[exo-17-2]{17.2}, \hyperref[exo-17-3]{17.3}, \hyperref[exo-17-9]{17.9}, \hyperref[exo-17-14]{17.14}\\[2mm]
L’axe en diagonale & \hyperref[exo-17-4]{17.4}, \hyperref[exo-17-5]{17.5}, \hyperref[exo-17-6]{17.6}, \hyperref[exo-17-7]{17.7}, \hyperref[exo-17-11]{17.11}\\[2mm]
Conservations et pavages & \hyperref[exo-17-10]{17.10}, \hyperref[exo-17-12]{17.12}, \hyperref[exo-17-15]{17.15}\\[2mm]
\bottomrule\end{tabularx}\par\vspace{4mm}
\otIntertitre{18. Solides : cubes, pavés et patrons}\begin{tabularx}{\linewidth}{@{}Yl@{}}\toprule
Descriptions des solides & \hyperref[exo-18-1]{18.1}, \hyperref[exo-18-2]{18.2}, \hyperref[exo-18-3]{18.3}, \hyperref[exo-18-4]{18.4}\\[2mm]
Patrons du cube & \hyperref[exo-18-5]{18.5}, \hyperref[exo-18-6]{18.6}, \hyperref[exo-18-8]{18.8}\\[2mm]
Patrons du pavé & \hyperref[exo-18-7]{18.7}, \hyperref[exo-18-9]{18.9}\\[2mm]
Assemblages et déplacements & \hyperref[exo-18-10]{18.10}, \hyperref[exo-18-11]{18.11}, \hyperref[exo-18-12]{18.12}, \hyperref[exo-18-13]{18.13}, \hyperref[exo-18-14]{18.14}, \hyperref[exo-18-15]{18.15}\\[2mm]
\bottomrule\end{tabularx}\par\vspace{4mm}
\clearpage
\otPageTitre{Repères et exercices}{Chapitres 19 à 21}\fontsize{10}{13}\selectfont
\otIntertitre{19. La proportionnalité en actes}\begin{tabularx}{\linewidth}{@{}Yl@{}}\toprule
Reconnaître et justifier & \hyperref[exo-19-1]{19.1}, \hyperref[exo-19-9]{19.9}, \hyperref[exo-19-15]{19.15}\\[2mm]
Fois plus, fois moins & \hyperref[exo-19-2]{19.2}, \hyperref[exo-19-3]{19.3}, \hyperref[exo-19-7]{19.7}\\[2mm]
Additionner des situations & \hyperref[exo-19-4]{19.4}, \hyperref[exo-19-5]{19.5}, \hyperref[exo-19-6]{19.6}, \hyperref[exo-19-11]{19.11}, \hyperref[exo-19-12]{19.12}\\[2mm]
Comparaisons et échelles & \hyperref[exo-19-8]{19.8}, \hyperref[exo-19-10]{19.10}, \hyperref[exo-19-13]{19.13}, \hyperref[exo-19-14]{19.14}\\[2mm]
\bottomrule\end{tabularx}\par\vspace{4mm}
\otIntertitre{20. Lire et croiser des données}\begin{tabularx}{\linewidth}{@{}Yl@{}}\toprule
Tableaux et totaux & \hyperref[exo-20-1]{20.1}, \hyperref[exo-20-2]{20.2}, \hyperref[exo-20-7]{20.7}, \hyperref[exo-20-12]{20.12}\\[2mm]
Barres et courbes & \hyperref[exo-20-3]{20.3}, \hyperref[exo-20-4]{20.4}, \hyperref[exo-20-6]{20.6}, \hyperref[exo-20-8]{20.8}\\[2mm]
Diagrammes circulaires & \hyperref[exo-20-5]{20.5}, \hyperref[exo-20-9]{20.9}\\[2mm]
Croiser et critiquer & \hyperref[exo-20-10]{20.10}, \hyperref[exo-20-11]{20.11}, \hyperref[exo-20-13]{20.13}, \hyperref[exo-20-14]{20.14}, \hyperref[exo-20-15]{20.15}\\[2mm]
\bottomrule\end{tabularx}\par\vspace{4mm}
\otIntertitre{21. Les probabilités : a chances sur b}\begin{tabularx}{\linewidth}{@{}Yl@{}}\toprule
a chances sur b & \hyperref[exo-21-1]{21.1}, \hyperref[exo-21-2]{21.2}, \hyperref[exo-21-3]{21.3}, \hyperref[exo-21-13]{21.13}\\[2mm]
Comparer par les issues & \hyperref[exo-21-4]{21.4}, \hyperref[exo-21-5]{21.5}, \hyperref[exo-21-10]{21.10}, \hyperref[exo-21-12]{21.12}, \hyperref[exo-21-14]{21.14}, \hyperref[exo-21-15]{21.15}\\[2mm]
Indépendance des répétitions & \hyperref[exo-21-6]{21.6}\\[2mm]
Arbres et tableaux à deux étapes & \hyperref[exo-21-7]{21.7}, \hyperref[exo-21-8]{21.8}, \hyperref[exo-21-9]{21.9}, \hyperref[exo-21-11]{21.11}\\[2mm]
\bottomrule\end{tabularx}\par\vspace{4mm}
\clearpage
